\section{Navigation for Source Section II}
\label{sec:theory-overview}

\sourceeqrange{Eqs. (2.1)--(2.53)}

\subsection*{Purpose}

Use this page as the dependency map for Section II.  It groups the source
objects by role: thermodynamic, variational, hydrodynamic, entropy production,
and wall transfer.

\subsection*{Derivation route}

\begin{itemize}
\item Separate equilibrium thermodynamics from nonequilibrium balance laws.
\item Track where the gradient entropy and gradient energy enter.
\item Derive the reversible stress and entropy-flux identities in components.
\item Record exactly which statements are local and which are integrated.
\end{itemize}

\begin{figure}[htbp]
\centering
\resizebox{\linewidth}{!}{\begin{tikzpicture}[
  font=\small,
  node distance=8mm and 10mm,
  box/.style={draw=black, line width=0.45pt, rounded corners=2pt,
    align=center, text width=32mm, minimum height=11mm, fill=gray!8},
  source/.style={draw=black, line width=0.45pt, rounded corners=2pt,
    align=center, text width=34mm, minimum height=11mm, fill=gray!16},
  open/.style={draw=black, dashed, line width=0.45pt, rounded corners=2pt,
    align=center, text width=38mm, minimum height=11mm, fill=white},
  arrow/.style={-{Latex[length=2.2mm]}, line width=0.5pt},
  darrow/.style={-{Latex[length=2.2mm]}, dashed, line width=0.5pt}
]
\node[source] (bulk) {II.A\\local van der Waals baseline};
\node[box, right=of bulk] (grad) {II.B\\gradient entropy and energy};
\node[box, right=of grad] (wall) {II.C\\equilibrium wall branch};
\node[box, below=of grad] (coeff) {\(M=CT+K\)\\coefficient map};
\node[box, right=of coeff] (stress) {Appendix A\\reversible stress};
\node[box, right=of stress] (hydro) {II.D\\local balance laws};
\node[open, below=of stress] (global) {Integrated energy/entropy requires heat, wall, and boundary transfer data};

\draw[arrow] (bulk) -- (grad);
\draw[arrow] (grad) -- (wall);
\draw[arrow] (grad) -- (coeff);
\draw[arrow] (coeff) -- (stress);
\draw[arrow] (stress) -- (hydro);
\draw[darrow] (wall) -- (global);
\draw[darrow] (hydro) -- (global);
\end{tikzpicture}
}
\caption{Dependency flow for the Section II derivations.  The diagram separates
local thermodynamic and hydrodynamic identities from the integrated
wall-transfer branch, which remains source-dependent.}
\label{fig:section-ii-dependency-flow}
\end{figure}

\subsection*{Verification route}

\begin{itemize}
\item Component residuals for gradient-force and stress-divergence identities.
\item Dimension checks for energy density, entropy density, stress, heat flux,
  entropy flux, and surface-energy terms.
\item Logical guardrails: local flux equality does not imply global equality
  without a boundary map.
\end{itemize}

\subsection*{Open issues}

Full global energy/entropy closure remains unresolved until wall transfer,
thermal boundary data, and admissible variations are source-mapped.
